% =============================================================================
% The PreDoc logit experiment, reproduced after four additions to the package.
%
% Every number in this document is transcribed from the executed notebook
% notebooks/Sampling/predoc_logit_v2.ipynb. Nothing is computed here and nothing
% is rounded beyond the digits the notebook printed.
%
% Every quantity attributed to the PreDoc is transcribed from
% _These/PreDoc/PreDoc_JeremyRieussec_2021_v2/main.tex, with the line given.
%
% Standalone. Own preamble, no bibliography, no \cite.
% =============================================================================
\documentclass[11pt,a4paper]{article}

\usepackage[utf8]{inputenc}
\usepackage[british]{babel}
\usepackage{amsmath, amssymb, amsthm}
\usepackage{booktabs}
\usepackage{longtable}
\usepackage{geometry}
\usepackage{float}
\usepackage{graphicx}
\usepackage{enumitem}
\usepackage{xcolor}
\usepackage[hidelinks]{hyperref}
\geometry{margin=2.2cm}

\providecommand{\norm}[1]{\left \lVert #1 \right \rVert}
\providecommand{\abs}[1]{\left \vert #1 \right \vert}
\newcommand{\VAI}{\textsf{VAI}}
\newcommand{\VAIVAC}{\textsf{VAI/VAC}}
\newcommand{\VAC}{\textsf{VAC}}
\newcommand{\NEW}{\textcolor{blue!60!black}{\textbf{new}}}

\title{The PreDoc logit experiment, reproduced\\
       \large four additions to \textsf{TrustRegionRadius.jl}, and what the runs say}
\author{}
\date{}

\begin{document}
\maketitle

\begin{abstract}
The variable-sampling experiment of the PreDoc needed four things the package did
not have: a multiple-alternative logit, a two-sided smoothing of the sample size,
the outer-product estimate of the paired variance, and the three sample management
schemes. All four are now in the package, and this report gives what was added,
why each was needed, how each was checked, and what the reproduction shows. The
experiment is the PreDoc's: five alternatives, $p = 10$, $N = 100\,000$ individuals,
$\lvert\mathcal{S}_0\rvert = 100$, $z_{0.95}$, and the trust-region constants
$(\eta_1,\eta_2,\gamma_1,\gamma_2,\gamma_3) = (0.01, 0.8, 0.5, 0.9, 1.5)$. The
central claim reproduces: the common-random-variable schemes converge in 26 and 28
iterations where independent draws do not converge in 300. Three qualifications
came out of the runs that the iteration counts hide. All three arms spend within
6\% of the same number of term evaluations, so the right statement is that one arm
converged on the budget the other spent without converging. The mechanism the
PreDoc offers for the effect is not confirmed by the natural diagnostic, which
ranks the schemes in the wrong order. And the overfitting hazard the PreDoc
attaches to \VAC{} is present and measurable, at 4 of 8 sample-size falls taken
with the true gradient rising, against 1 of 6 under \VAIVAC{}.
\end{abstract}

\tableofcontents

% =============================================================================
\section{What this supersedes}
% =============================================================================

An earlier report, \texttt{report/task14-predoc-logit/}, recorded the same
experiment attempted under the constraint that nothing in the package be modified.
It reported four departures from the PreDoc as gaps and reported one comparison as
refuted rather than made. That report is a correct record of what was possible then
and is superseded by this one on every point where the two disagree. The four gaps
are now features; the refuted comparison is now made directly, in
Section~\ref{sec:variance}.

% =============================================================================
\section{The experiment, reformulated}
% =============================================================================

\subsection{The problem}

Individual $n$ chooses among $J$ alternatives. Alternative $j$ carries an attribute
vector $x_{nj} \in \mathbb{R}^p$ and a representative utility linear in the
parameters, $V_{nj} = \beta^\top x_{nj}$. Under independent Gumbel errors the
choice probabilities take the closed form

\begin{equation}
  P_{nj} \;=\; \frac{e^{\beta^\top x_{nj}}}{\sum_{l} e^{\beta^\top x_{nl}}},
  \label{eq:logit}
\end{equation}

and maximum likelihood over a population of $N$ individuals is the minimisation of

\begin{equation}
  f(\beta) \;=\; -\frac{1}{N}\sum_{n=1}^{N}\sum_{j=1}^{J} y_{nj}\ln P_{nj},
  \qquad y_{nj} = 1 \text{ iff } n \text{ chose } j.
  \label{eq:objective}
\end{equation}

The PreDoc takes $J = 5$, $p = 10$ and $N = 100\,000$, on data generated from the
model itself so that the information identity holds and the outer-product Hessian
is justified.

\subsection{The method}

The method is a trust-region algorithm on a Monte-Carlo model
$m_k(s) = \tilde f_k + \tilde g_k^\top s + \tfrac12 s^\top \tilde B_k s$, with
three additions to the standard loop.

\paragraph{The sample size candidate.} From a confidence statement about the
achieved decrease, the PreDoc's Step 5 asks for

\begin{equation}
  N_k^{+} \;=\; \left\lceil
      \frac{\sigma_k^2 \bigl(\Phi^{-1}(1-\beta_k)\bigr)^2}{\tilde\mu_k^2}
  \right\rceil,
  \qquad \tilde\mu_k = \tilde f_k(x_k + s_k) - \tilde f_k(x_k).
  \label{eq:step5}
\end{equation}

\paragraph{The sample update function.} Step 6 passes that candidate through
$\Psi_k$ before it becomes $N_{k+1}$. The \emph{naive smoothing}, written
$\mathrm{Naive}(a,b)$, is the two-sided bound

\begin{equation}
  a N_k \;\le\; N_{k+1} \;\le\; b N_k, \qquad 0 < a \le 1,\; b > 1,
  \label{eq:naive}
\end{equation}

alongside a floor that rises by a fixed increment,
$N^{\min}_{k+1} = N^{\min}_{k} + N_{incr}$.

\paragraph{The sample management scheme.} How $\mathcal{N}_k$ becomes
$\mathcal{N}_{k+1}$ is a separate choice from how large it is, and the PreDoc gives
three answers. The names differ between its method section and its figures, and
both are kept here because both appear in the source.

\begin{center}
\begin{tabular}{llp{8.4cm}}
\toprule
figures & method & what happens between iterations \\
\midrule
\VAI{}    & IRV   & redraw independently of the previous batch \\
\VAIVAC{} & I/CRV & on a rise, keep the batch and append fresh draws,
                    $\mathcal{N}_{k+1} = \mathcal{N}_k \cup \{\xi_{N_k+1},\dots\}$;
                    on a fall, keep a uniform subsample,
                    $\mathcal{N}_{k+1} \subset \mathcal{N}_k$ \\
\VAC{}    & CRV   & the first $N_k$ individuals,
                    $\mathcal{N}_k = \{\xi_1,\dots,\xi_{N_k}\}$; finite population only \\
\bottomrule
\end{tabular}
\end{center}

\paragraph{Which variance.} Equation~\eqref{eq:step5} needs $\sigma_k^2$, and the
PreDoc offers two, labelling its runs \emph{tv} and \emph{sHs}. The first is the
empirical variance of the individual decreases. The second comes from a first-order
expansion of $F(\cdot,\xi)$ about $x_k$, under which
$\mathrm{Var}[F(x_k+s_k,\xi) - F(x_k,\xi)] \approx
 s_k^\top \mathrm{Var}[\nabla F(x_k,\xi)]\, s_k$, and replacing
$\mathrm{Var}[\nabla F]$ by $\tilde B_k - \tilde g_k \tilde g_k^\top$ gives

\begin{equation}
  \hat\sigma_{N_k}^2 \;\approx\; s_k^\top \tilde B_k s_k - \bigl(\tilde g_k^\top s_k\bigr)^2 .
  \label{eq:opvar}
\end{equation}

Its appeal is that both terms are already formed: $s_k^\top\tilde B_ks_k$ is needed
for $\rho_k$ and $\tilde g_k$ is in memory for the model.

% =============================================================================
\section{Audit: what the package already had}
% =============================================================================

Five of the eight components needed no work.

\begin{center}
\begin{tabular}{lll}
\toprule
PreDoc & package & note \\
\midrule
Step 5, equation~\eqref{eq:step5} & \texttt{CertifiedDecrease} & on paired differences \\
Step 3, truncated CG & \texttt{SteihaugCG} & \\
Step 7, equation for $\Delta_{k+1}$ & \texttt{RDelta} & $\gamma_1,\gamma_2,\gamma_3$ as given \\
$\tilde B_k$, the OP approximation & \texttt{BHHHModel} & \\
a subsampled true Hessian & \texttt{ExactHessian} & the PreDoc's \emph{TH} \\
\bottomrule
\end{tabular}
\end{center}

\texttt{CertifiedDecrease} is worth a word because the correspondence is exact and
not obvious. It sizes the batch from the sample mean and variance of the paired
differences $D_i = F(x_k,\xi_i) - F(x_k+s_k,\xi_i)$, taken under common random
numbers, which is \eqref{eq:step5} with the sign of $\tilde\mu_k$ flipped so the
mean decrease is positive. Pairing is not an optimisation: with independent draws
at the two ends of the step, $\mathrm{Var}(D_i)$ does not shrink as
$\norm{s_k} \to 0$, and the rule becomes unusable near a solution.

Three thresholds also need care. The package separates acceptance from radius
scaling: acceptance is decided by $\eta$ and scaling by $\eta_1$ and $\eta_2$,
whereas the PreDoc accepts on its own $\eta_1$. Setting $\eta = \eta_1 = 0.01$
reproduces it.

The confidence level is passed as $p$ and the quantile derived from it. The
notebook checks that this reaches the PreDoc's constant rather than assuming it:
\texttt{CertifiedDecrease(p = 0.95).z} is $1.6448534923$ against the quoted
$z_{0.95} \approx 1.6449$, a difference of $4.651\times 10^{-5}$, which is agreement
to the four digits quoted.

% =============================================================================
\section{The four additions}
% =============================================================================

\subsection{\texorpdfstring{\texttt{MultinomialLogit}}{MultinomialLogit}}

A \texttt{LikelihoodProblem} carrying the $M \times J \times p$ attribute array and
the chosen alternative. It implements \eqref{eq:logit} and \eqref{eq:objective}
directly. Writing $\bar x_n = \sum_j P_{nj} x_{nj}$ for the probability-weighted
mean attribute, the per-observation score and Hessian are

\begin{equation}
  s_n = \bar x_n - x_{n c_n},
  \qquad
  H_n = \sum_j P_{nj}\, x_{nj} x_{nj}^\top - \bar x_n \bar x_n^\top ,
  \label{eq:mnl-derivs}
\end{equation}

so $H_n$ is the covariance of the attribute vector under the choice probabilities.
It is positive semidefinite, and the objective is therefore convex. This matters
downstream: \texttt{BHHHModel} is positive semidefinite by construction and so can
never report negative curvature, and on a convex problem that costs nothing. The
comparison against the exact Hessian in Section~\ref{sec:models} is therefore about
approximation quality alone.

The default population is the PreDoc's $100\,000 \times 5 \times 10$, which is 40 MB.

\paragraph{The generating noise, and a point we do not settle.} The PreDoc says
that during generation ``some noise is injected in the data by adding some
perturbation, with mean zero, to the true model parameters $\beta^*$'', and
concludes that the information identity holds. Read literally, a separate
$\beta_n = \beta^* + \varepsilon_n$ per individual is a \emph{mixed} logit
generating process fitted with a plain logit, which is misspecification, and the
identity then fails at every $\beta$ and does not improve with $M$.

The keyword \texttt{noise} implements the literal reading, and measurement settles
what it costs: at $\beta^*$ with $M = 4000$, the BHHH identity error is $0.064473$
at \texttt{noise = 0} and $1.365955$ at \texttt{noise = 1.5}. The runs in this
report all use \texttt{noise = 0}, the correctly specified case, so that the premise
Table~\ref{tab:identity} checks is the premise the runs rely on. We have not
reconciled the two readings and do not claim to.

\subsection{\texorpdfstring{\texttt{SmoothedSize}}{SmoothedSize}}

A wrapper over any \texttt{SamplingRule}, following the package's own idiom for
composition. The inner rule proposes $N^{+}_{k+1}$ and the wrapper reports

\begin{equation}
  N_{k+1} = \min\bigl\{ b N_k,\; \max\{ a N_k,\; N^{\min}_k,\; N^{+}_{k+1}\} \bigr\},
  \qquad N^{\min}_{k+1} = N^{\min}_k + N_{incr},
\end{equation}

which is \eqref{eq:naive} together with the PreDoc's floor. Every trait is
forwarded, so the wrapper is transparent to the problem-compatibility checks, to
the oracle's cap logic and to the solver.

Two design points are worth recording. The gradient and objective channels are
smoothed separately against their own previous sizes, so a rule that distinguishes
them keeps the distinction; a rule that returns one size for both, which is every
rule that pairs, is unaffected. And when the rising floor exceeds the ceiling
$bN_k$, which happens most easily just after a sharp fall, the floor wins: it is a
requirement the caller stated outright, where the ceiling is a smoothing
preference.

This subsumes what was there before. \texttt{CertifiedDecrease(monotone = true)} is
the special case $a = 1$, and \texttt{monotone = false} is $a = 0$, $b = \infty$.
Neither is the PreDoc's setting, and Table~\ref{tab:all} shows the two extremes
behaving quite differently from $\mathrm{Naive}(0.75, 2)$.

\subsection{\texorpdfstring{\texttt{CertifiedDecrease(variance = ...)}}{CertifiedDecrease(variance)}}
\label{sec:variance-feature}

A keyword taking \texttt{:empirical} (the default, unchanged) or
\texttt{:outer\_product}, together with a trait \texttt{paired\_variance\_kind} that
the solver reads to decide what to hand the rule. The outer-product estimate is
computed by \texttt{paired\_op\_variance}.

It is formed as the sample variance of the scalars $\nabla F_i^\top s_k$ rather than
as the difference $s^\top \tilde B s - (\tilde g^\top s)^2$ of two quadratic forms.
The two are algebraically the same, and the first cannot come out negative when the
two terms nearly cancel. This is not a hypothetical concern: it is exactly the
failure the earlier report had to record when it tried to reconstruct the quantity
after the fact from an archived trace.

Both estimates are now recorded on the trace whenever the problem supplies scores,
under

\begin{center}
\texttt{:paired\_variance\_trajectory} \qquad and \qquad
\texttt{:paired\_op\_variance\_trajectory},
\end{center}

so the estimate a run did not use is auditable against the one it did, on the same
batch at the same iterate. The cost is one score matrix over the objective batch per
paired iteration.

One distinction is easy to lose and worth stating. The outer product appears here as
an estimate of $\mathrm{Var}[\nabla F]$, not as an approximation to the Hessian. It
does \emph{not} require the information identity, and it remains legitimate on a
misspecified model where BHHH as a Hessian is not.

\subsection{\texorpdfstring{\texttt{FiniteSumNLP(scheme = ...)}}{FiniteSumNLP(scheme)}}

A field on both sampled oracles, taking \texttt{:independent}, \texttt{:nested} or
\texttt{:prefix}, implementing IRV, I/CRV and CRV to the definitions quoted above.
Within an iteration the batch is fixed and every evaluation shares it, which is
untouched; the scheme decides only what happens between iterations.

\texttt{:prefix} needs a finite population and is rejected on an expectation, which
is the PreDoc's own restriction of CRV. It consumes no random numbers at all, which
the notebook checks by comparing the oracle's generator state across five
resamples. \texttt{:nested} is supported on both classes through two small
functions the problem type can extend, with methods for index vectors and for
Gaussian draws.

% =============================================================================
\section{How the four were checked}
% =============================================================================

Forty-five checks, each against the definition rather than against the
implementation. The full suite of $1452$ pre-existing package tests also passes
unchanged, which is what establishes that the defaults are untouched.

\begin{center}
\small
\setlength{\tabcolsep}{4pt}
\begin{tabular}{lp{8.6cm}r}
\toprule
feature & check & result \\
\midrule
\texttt{MultinomialLogit}
  & gradient against central differences of the objective & rel.\ $5.05\times10^{-10}$ \\
  & Hessian against central differences of the gradient   & rel.\ $5.74\times10^{-10}$ \\
  & Hessian positive semidefinite                          & $\lambda_{\min} = 0.670$ \\
  & $f(0) = \ln J$ over the whole population               & $6.66\times10^{-16}$ \\
  & mean of the scores equals the gradient                 & $0$ \\
\midrule
\texttt{SmoothedSize}
  & every transition of a 300-iteration run inside \eqref{eq:naive} & $0$ violations \\
  & the same run's floor $10 + 10k$ respected               & $0$ violations \\
  & the bare rule does leave those bounds                   & ratios in $[0.10, 7.87]$ \\
\midrule
variance
  & \texttt{paired\_op\_variance} equals $s^\top\tilde Bs - (\tilde g^\top s)^2$ & rel.\ $<10^{-9}$ \\
  & the two estimates agree as $\norm{s}\to0$               & ratio $1.0059 \to 1.0000$ \\
  & rejected on a problem with no scores                    & \texttt{ArgumentError} \\
\midrule
schemes
  & \texttt{:prefix} batch is $1{:}N_k$ at every iteration  & holds \\
  & \texttt{:prefix} consumes no randomness                 & generator unchanged \\
  & \texttt{:nested} keeps the batch on a rise              & $\mathcal{N}_k \subseteq \mathcal{N}_{k+1}$ \\
  & \texttt{:nested} subsets it on a fall                   & $\mathcal{N}_{k+1} \subset \mathcal{N}_k$ \\
  & \texttt{:independent} does \emph{not} nest              & confirmed \\
\bottomrule
\end{tabular}
\end{center}

The last of those matters as much as the others: without it the nesting checks
would pass on a scheme that nested by accident, and would be evidence of nothing.

% =============================================================================
\section{Results}
% =============================================================================

Every arm below carries the PreDoc's trust-region constants, $x_0 = 0$,
$\Delta_0 = 0.1\norm{\tilde g_0}$ with $\tilde g_0$ over the initial batch of
$\lvert\mathcal{S}_0\rvert = 100$, and the same seed. Arms differ only in what they
are named for. Two settings are ours and not the PreDoc's, which states neither: a
cap of 300 iterations and a gradient tolerance of $10^{-6}$.

Every run is scored after the fact by the same code, on the true objective and true
gradient over all $100\,000$ individuals at the point it stopped, never on the
solver's own stopping test. An arm reporting \texttt{max\_iter} is reported as not
having converged, not as having taken that many iterations.

\subsection{The premise: does the information identity hold?}

\begin{table}[H]
\centering
\caption{The information identity, measured over all $100\,000$ individuals.
$B_{err}$ is the BHHH relative error, $W_{err}$ the BHHH-2 one, and the gap between
them is the rank-one term that vanishes at a stationary point.}
\label{tab:identity}
\begin{tabular}{lrrrr}
\toprule
point & $B_{err}$ & $W_{err}$ & $BW_{gap}$ & $\norm{g}$ \\
\midrule
$\beta^*$ (the truth)     & 0.011950 & 0.011949 & 0.000010 & $3.5033\times10^{-3}$ \\
$\beta^* + 0.5$           & 0.832126 & 0.599609 & 0.342115 & $6.7378\times10^{-1}$ \\
$0$ (the PreDoc's $x_0$)  & 0.245365 & 0.157656 & 0.402691 & $1.0094\times10^{0}$ \\
\bottomrule
\end{tabular}
\end{table}

The identity holds at $\beta^*$ to $1.2\%$, and the residual is the
$O(M^{-1/2})$ finite-sample discrepancy rather than a modelling error: in the
population the two sides are identically equal. It degrades to $83\%$ half a unit
away, which is the behaviour BHHH is known for and the reason the trust region is
the safeguard.

\subsection{The four questions}

\begin{table}[H]
\centering
\caption{Every arm, scored identically. $f$ and $\norm{\nabla f}$ are over the whole
population at the stopping point. \emph{terms} is the cumulative count of individual
evaluations, which is the cost; iterations are not. \texttt{max\_it.} and
\texttt{first\_ord.} abbreviate the statuses \texttt{max\_iter} and
\texttt{first\_order}.}
\label{tab:all}
\footnotesize
\setlength{\tabcolsep}{3pt}
\begin{tabular}{lrlrrrrrr}
\toprule
arm & iters & status & $N_{\min}$ & $N_{\max}$ & terms & $f$ (true) & $\norm{\nabla f}$ & acc \\
\midrule
raw, \VAI{}                 & 300 & \texttt{max\_it.}  & 10  & 3946    & $1.982\times10^{4}$ & 1.510823 & $3.284\times10^{-1}$ & 0.6509 \\
Naive(0.75,2), \VAI{}       & 300 & \texttt{max\_it.}  & 50  & 4717    & $9.794\times10^{5}$ & 0.782861 & $2.683\times10^{-2}$ & 0.6924 \\
monotone ($a=1$), \VAI{}    & 120 & \texttt{first\_ord.} & 100 & 100000  & $1.725\times10^{6}$ & 0.781998 & $9.002\times10^{-8}$ & 0.6915 \\
Naive, \VAIVAC{}            &  26 & \texttt{first\_ord.} & 50  & 100000  & $9.627\times10^{5}$ & 0.781998 & $7.107\times10^{-8}$ & 0.6915 \\
Naive, \VAC{}               &  28 & \texttt{first\_ord.} & 50  & 100000  & $1.019\times10^{6}$ & 0.781998 & $2.926\times10^{-8}$ & 0.6915 \\
Naive \emph{tv}, \VAIVAC{}  &  26 & \texttt{first\_ord.} & 50  & 100000  & $9.627\times10^{5}$ & 0.781998 & $7.107\times10^{-8}$ & 0.6915 \\
Naive \emph{sHs}, \VAIVAC{} &  24 & \texttt{first\_ord.} & 50  & 100000  & $8.304\times10^{5}$ & 0.781998 & $4.736\times10^{-7}$ & 0.6915 \\
Naive, \VAIVAC{}, exact $H$ &  31 & \texttt{first\_ord.} & 50  & 100000  & $1.064\times10^{6}$ & 0.781998 & $1.210\times10^{-9}$ & 0.6915 \\
\bottomrule
\end{tabular}
\end{table}

\paragraph{1. Is the raw certified size unstable?} Yes, and more sharply than
``unstable'' conveys. The size sits at $N_{\min} = 10$ for $79\%$ of the run, and
the transition ratio $N_{k+1}/N_k$ ranges over $[0.0025, 263.07]$, with 9 of 299
transitions jumping by more than fourfold. The cause is arithmetic and not
numerical: \eqref{eq:step5} reads $\hat\sigma$ and $\hat\delta$ off the \emph{same}
step, and when the decrease nearly vanishes the requirement jumps by orders of
magnitude with nothing to damp it.

The consequence is worse than an ugly plot. With ten individuals and ten
parameters the batch problem is degenerate, so the run is not solving the problem:
it ends at $f = 1.510823$ against $\ln 5 = 1.609438$ at the starting point, having
moved almost nowhere in $300$ iterations. It is also by far the cheapest arm, at
$1.982\times10^{4}$ terms, which is what a run that does nothing costs.

\paragraph{2. Does the naive smoothing fix it?} It improves the run by a long way
and does not fix it. $\mathrm{Naive}(0.75,2)$ with the floor brings the true
objective from $1.510823$ to $0.782861$ and the true gradient from $0.3284$ to
$0.02683$, and the bound holds at all $299$ transitions. But the arm still reports
\texttt{max\_iter}.


\begin{figure}[H]
\centering
\includegraphics[width=0.86\textwidth]{figs/predoc_mnl_fig2_smoothing.pdf}
\caption{The three sample update functions under independent draws, on a
logarithmic scale, with the rising floor $N^{\min}_k = 10 + 10k$ shown for
reference. The raw rule sits on $N_{\min}$ and spikes; $\mathrm{Naive}(0.75,2)$
tracks the floor; $a = 1$ climbs to the population and stays.}
\label{fig:smoothing}
\end{figure}

What converges under independent draws is the extreme setting $a = 1$: 120
iterations, and $1.725\times10^{6}$ terms, which is the most expensive arm in the
table by $60\%$. So the lower bound is doing real work here, and $a = 0.75$ is not
tight enough on its own on this problem. That is a qualification of the PreDoc's
report that $\mathrm{Naive}(0.75,2.0)$ ``empirically shows some better results'',
not a contradiction of it: the PreDoc compares Naive against CMS-cum, and the
monotone arm is neither.

\paragraph{3. Do the common-variable schemes converge in fewer iterations?} Yes,
and the gap is not marginal: 26 iterations under \VAIVAC{} and 28 under \VAC{},
against an independent arm that does not converge in 300. This is the PreDoc's
central empirical claim about sample management and it reproduces.


\begin{figure}[H]
\centering
\includegraphics[width=0.86\textwidth]{figs/predoc_mnl_fig3_schemes.pdf}
\caption{The same rule and the same seed under the three sample management
schemes. The two common-variable arms reach the population and stop; the
independent arm runs to the iteration cap without doing so.}
\label{fig:schemes}
\end{figure}

Three qualifications follow from the same table.

First, the iteration count is not the cost. The three arms spend
$9.794\times10^{5}$, $9.627\times10^{5}$ and $1.019\times10^{6}$ terms, all within
$6\%$ of each other, because the converging runs reach the full population and their
late iterations are correspondingly expensive. The defensible statement is not that
common variables are cheaper. It is that on the same budget one arm converged and
the other did not, which is a stronger claim and a different one.

Second, the mechanism is not confirmed. The PreDoc attributes the effect to common
variables stopping $\tilde f$ from moving for reasons unrelated to the step. The
natural diagnostic, the median of $\abs{\tilde f_{k+1} - \tilde f_k}$ against the
median achieved decrease over a common window, gives $1.430$, $4.267$ and $6.337$
for \VAI{}, \VAIVAC{} and \VAC{}, which is the wrong order. The absolute movement
does fall as predicted, from $7.605\times10^{-2}$ to $6.796\times10^{-3}$ and
$8.638\times10^{-3}$, but so does the decrease, and by more. The comparison is
confounded in any case: at a matched iteration index the \VAI{} run is still far out
taking large steps while the \VAIVAC{} run is nearly converged. We report the
diagnostic and do not draw the conclusion from it.

Third, all three arms reach the same objective to six figures, $0.781998$ against
$0.782861$, and the same accuracy to three. What separates them is whether they got
there, not where they got.

\paragraph{The \VAC{} hazard.} The PreDoc warns that the prefix scheme can overfit
its own prefix: while $N_k$ is small the run minimises a different function, whose
minimiser is some distance from the one sought, and the symptom is a spike then a
fall in the sample size, possibly cycling.

A fall on its own is not the hazard, since with $a = 0.75$ the rule is permitted to
fall and should whenever the decrease is large. The hazard is a fall taken while the
\emph{true} gradient is rising, and that is on the trace.

\begin{table}[H]
\centering
\caption{Falls in the sample size, and whether the run was going backwards across
them. The comparison that carries weight is the last two rows, which converged
under the same rule and seed.}
\label{tab:hazard}
\begin{tabular}{lrrrr}
\toprule
arm & iters & falls & falls with $\norm{\nabla f}$ up & worst rise \\
\midrule
Naive(0.75,2), \VAI{} & 300 & 56 & 33 & $3.316\times$ \\
Naive, \VAIVAC{}      &  26 &  6 &  1 & $1.149\times$ \\
Naive, \VAC{}         &  28 &  8 &  4 & $1.779\times$ \\
\bottomrule
\end{tabular}
\end{table}

The signature is present, and it is proportionally worse under \VAC{} than under
\VAIVAC{}, at 4 of 8 against 1 of 6 and a worst rise of $1.78\times$ against
$1.15\times$. That is the ordering the PreDoc predicts, on runs identical in every
other respect, and it is the mechanism it describes: \VAC{} reuses one prefix while
\VAIVAC{} refreshes it at the margin.

The \VAI{} row has the largest raw count, and it is not evidence about sample
management: a run that never converged has a gradient moving for its own reasons.

The hazard did not prevent convergence here, and it should not be expected to. The
$100\,000$ individuals are generated independently, so the first hundred are an
ordinary random sample. Sorted or clustered data is where it would bite, and this
reproduction does not construct that case.

\subsection{The two variance estimates}
\label{sec:variance}

The earlier report could not make this comparison. Reconstructing
$s^\top\tilde Bs$ after the fact from an archived trace was refuted by its own
check, on $84$ of $300$ iterations, for two independent reasons: the achieved
reduction is not the difference of successive objective values on a rejected step,
and under sampling the two values sit on different batches. Both estimates are now
formed inside the iteration, on one batch at one iterate, and recorded.

\begin{table}[H]
\centering
\caption{The two estimates against each other, over the 26 paired iterations of the
\VAIVAC{} run.}
\label{tab:variance}
\begin{tabular}{lr}
\toprule
iterations with both finite & 26 \\
median ratio $\hat\sigma^2_{sHs} / \hat\sigma^2_{tv}$ & 1.003150 \\
ratio range & $[0.999912,\; 1.218152]$ \\
iterations with ratio outside $[0.5, 2]$ & 0 \\
negative outer-product values & 0 of 26 \\
\bottomrule
\end{tabular}
\end{table}


\begin{figure}[H]
\centering
\includegraphics[width=0.86\textwidth]{figs/predoc_mnl_fig5_variances.pdf}
\caption{The empirical and outer-product estimates of $\mathrm{Var}(D_i)$, formed
on the same batch at the same iterate of the \VAIVAC{} run, with $\norm{s_k}$ on
the same axes. The two are indistinguishable at this scale over most of the run,
which is Table~\ref{tab:variance} seen as a trajectory.}
\label{fig:variances}
\end{figure}

The agreement is close: the median discrepancy is $0.3\%$ and the worst is $22\%$.
As a rule the choice makes little difference and slightly favours \emph{sHs}, which
converges in 24 iterations rather than 26 and at $8.304\times10^{5}$ terms rather
than $9.627\times10^{5}$.

Two cautions. The approximation is first order in $\norm{s_k}$, so the spread in
the ratio is the step length and not an estimation error; the notebook confirms the
direction directly, with the ratio running $1.0059 \to 1.0008 \to 1.0001 \to 1.0000$
as $\norm{s}$ falls through four decades. And the close agreement is a property of
this problem, where Table~\ref{tab:identity} shows the identity holding to $1.2\%$.
It should not be carried to a model where it does not.

\subsection{BHHH against the exact Hessian}
\label{sec:models}

\begin{table}[H]
\centering
\caption{The two model Hessians under the nested scheme, everything else fixed.}
\begin{tabular}{lrrrr}
\toprule
model & iters & terms & $\norm{\nabla f}$ (true) & CG iterations \\
\midrule
\texttt{BHHHModel}    & 26 & $9.627\times10^{5}$ & $7.107\times10^{-8}$ & 55 \\
\texttt{ExactHessian} & 31 & $1.064\times10^{6}$ & $1.210\times10^{-9}$ & 53 \\
\bottomrule
\end{tabular}
\end{table}

BHHH reaches the same objective in fewer iterations and at lower cost, and the exact
Hessian reaches a smaller final gradient. Neither is a surprise at $p = 10$ with the
identity holding to $1.2\%$: the two models are close, and the difference is within
what one seed shows.

The cost argument is about $p$ and not about this table. The exact Hessian costs
$O(NJp^2)$ per batch against $O(NJp)$ for the score matrix, and the score matrix is
computed for the gradient anyway, so the outer product is obtained free. At $p = 10$
the constant is small enough that the two run in comparable time.

One reporting note. \texttt{neval\_hprod} reads $0$ for the BHHH arm, because
\texttt{OuterProductOperator} forms $Bv$ from the score matrix without going through
the model's \texttt{hprod!}, so the counter never sees it. It did one product per CG
iteration. The CG line is the one to read when the model is an outer product.

% =============================================================================
\section{What is not reproduced}
% =============================================================================

Four things, stated so that the reproduction is not read as more complete than it is.

\begin{enumerate}[leftmargin=*]
\item \textbf{CMS and CMS-cum smoothing.} The PreDoc compares $\mathrm{Naive}$
      against $\mathrm{CMS\text{-}cum}$, and only $\mathrm{Naive}$ is implemented.
      CMS holds the size constant for $m = \lceil \epsilon^\beta_k/\tilde\mu_k\rceil$
      iterations, which needs the rule to carry a countdown across iterations rather
      than answer from the current state. That is a different shape of rule, not a
      different constant, and \texttt{SmoothedSize} is not it.

\item \textbf{A separate Hessian sample $\mathcal{H}_k$.} The PreDoc's Step 1 draws
      $\mathcal{H}_k$, possibly non-uniformly, and then assumes
      $\mathcal{H}_k \subset \mathcal{S}_k$ throughout. Here the model is built on
      the gradient batch, which is the case $\mathcal{H}_k = \mathcal{S}_k$.

\item \textbf{The generating noise.} Discussed above: the keyword exists and the
      runs use $\texttt{noise} = 0$. The two readings of the PreDoc's sentence are
      not reconciled.

\item \textbf{A convergence criterion.} The PreDoc states none, and its iteration
      counts are read off figures. The cap of $300$ and the tolerance of $10^{-6}$
      are ours. This is why comparisons against arms reporting \texttt{max\_iter}
      are phrased as failures to converge within the cap.
\end{enumerate}

% =============================================================================
\section{Where things are}
% =============================================================================

\begin{center}
\begin{tabular}{ll}
\toprule
notebook & \texttt{notebooks/Sampling/predoc\_logit\_v2.ipynb} \\
driver & \texttt{\_audit/nb7\_predoc\_logit\_mnl.jl} \\
\midrule
\texttt{MultinomialLogit} & \texttt{src/Likelihood/models.jl} \\
\texttt{SmoothedSize}, \texttt{paired\_variance\_kind} & \texttt{src/Sampling/rules.jl} \\
schemes, \texttt{paired\_op\_variance} & \texttt{src/Sampling/oracles.jl} \\
recording both variances & \texttt{src/Trust-region/common.jl} \\
\midrule
superseded report & \texttt{report/task14-predoc-logit/} \\
\bottomrule
\end{tabular}
\end{center}

The notebook is six figures and eight tables, executed end to end with sequential
execution counts and no errors. Nothing in this report is computed here; every
number is transcribed from that execution.

\end{document}
